\begin{remark}\label{fga-rem-9-2-11}\dcref{9.2.11}\source{fga-explained:page-0266}\uses{fga-thm-9-2-5}
There is a derived proof of the comparison with the Picard functor.  For a ringed
space $T$, $\operatorname{Pic}(T)=H^1(T,\mathcal O_T^\times)$, and the Leray
spectral sequence for $f_T:X_T\to T$ gives
\[
0\to H^1(T,\mathcal O_T^\times)\to H^1(X_T,\mathcal O_{X_T}^\times)
\to H^0(T,R^1f_{T*}\mathcal O_{X_T}^\times)\to H^2(T,\mathcal O_T^\times).
\]
If $\mathcal O_T\to f_{T*}\mathcal O_{X_T}$ is universally an isomorphism,
the first map identifies the relative Picard functor with its Zariski, etale,
or fppf sheafification.  If $f$ has a section, the last map is injective and
the comparison is an isomorphism.  In the etale case the final obstruction is
the cohomological Brauer class in $\operatorname{Br}'(T)$.
\end{remark}

\begin{remark}\label{fga-rem-9-3-9}\dcref{9.3.9}\source{fga-explained:page-0270}\uses{}
The map from ordered effective divisors factors through the symmetric product:
the addition map $X^d\to\operatorname{Div}_{X/S}$ is invariant under the
permutation action and induces the usual map $\operatorname{Sym}^d(X/S)\to
\operatorname{Div}_{X/S}$.  For a proper morphism $f:X\to S$ and a coherent
flat sheaf $\mathcal F$, there is a coherent $\mathcal O_S$-module $Q$ with
\[
\operatorname{Hom}_S(Q,\mathcal M)\simeq
f_*(\mathcal H om(\mathcal F,\mathcal F\otimes f^*\mathcal M))
\]
functorially in quasi-coherent $\mathcal M$.  At $s\in S$, the following are
equivalent: $Q_s$ is free; the functor $\mathcal M\mapsto
f_*(\mathcal F\otimes f^*\mathcal M)$ is right exact; base change is an
isomorphism for all $\mathcal M$; and
$H^0(X,\mathcal F)\otimes k(s)\to H^0(X_s,\mathcal F_s)$ is surjective.
They are implied by $H^1(X_s,\mathcal F_s)=0$.
\end{remark}

\begin{remark}\label{fga-rem-9-4-18}\dcref{9.4.18}\source{fga-explained:page-0279}\uses{fga-thm-9-4-8}
The principal refinements of the existence theorem for the Picard scheme are
Mumford's, Grothendieck's, and the later representability results for the local
and non-flat Picard functors.  Each replaces the basic hypotheses by a new
boundedness or algebraization argument.
\end{remark}

\begin{remark}\label{fga-rem-9-4-19}\dcref{9.4.19}\source{fga-explained:page-0282}\uses{}
The local Picard scheme treats the punctured spectrum of a noetherian local
algebra, while the non-flat theory treats proper schemes over complete local
rings.  Both enter the proof of Samuel's conjecture: a complete local UFD over
an algebraically closed field has $A[[t]]$ again a UFD.  For a resolution of a
normal surface singularity, the local Picard group is obtained from the Picard
group of the resolution by quotienting by the classes of exceptional components;
formal geometry identifies the latter with the inverse limit of the Picard
groups of the infinitesimal thickenings.
\end{remark}

\begin{remark}\label{fga-rem-9-5-2}\dcref{9.5.2}\source{fga-explained:page-0286}\uses{}
For a finite-type group scheme $G$ over a field, the reduction $G_{\mathrm{red}}$
need not be a subgroup scheme: the product $G_{\mathrm{red}}\times G_{\mathrm{red}}$
may fail to be reduced.  Additional hypotheses are required for the reduced
subscheme to inherit the group law.
\end{remark}

\begin{remark}\label{fga-rem-9-5-6}\dcref{9.5.6}\source{fga-explained:page-0287}\uses{fga-cor-9-4-18-3,fga-thm-9-5-4}
Theorem~9.5.4 remains valid for every proper $X/k$, without assuming that $X$
is integral.  In the proof, the general existence theorem replaces the integral
case used for the quasi-projectivity argument.
\end{remark}

\begin{remark}\label{fga-rem-9-5-8}\dcref{9.5.8}\source{fga-explained:page-0288}\uses{fga-thm-9-5-4}
The projectivity of the identity component of the Picard scheme admits several
proofs.  One inducts on the dimension using a general hyperplane section and
proves that the induced pullback on Picard varieties has finite kernel.  In
characteristic zero, the kernel vanishes by the argument with unramified cyclic
covers and connectedness of ample divisors; in arbitrary characteristic it is
controlled by the corresponding finite unipotent group.
\end{remark}

\begin{proposition}\label{fga-pro-9-5-10}\dcref{9.5.10}\source{fga-explained:page-0290}\uses{fga-thm-9-4-8}
Let $X/k$ be proper with $k$ a field, and assume the fppf Picard functor is
represented by $\operatorname{Pic}_{X/k}$.  An invertible sheaf $\mathcal L$ on
$X$ is algebraically equivalent to $\mathcal O_X$ if and only if the point of
$\operatorname{Pic}_{X/k}$ represented by $\mathcal L$ lies in
$\operatorname{Pic}^0_{X/k}$.
\end{proposition}
\begin{proof}
An algebraic family of line bundles over a connected base maps into one
connected component, and a family joining $\mathcal O_X$ to $\mathcal L$
therefore gives a point of $\operatorname{Pic}^0$.  Conversely, pull back the
universal line bundle to an irreducible curve in a connected fppf cover of
$\operatorname{Pic}^0$ through the given point.  The two fibres over points of
the curve are algebraically equivalent, and one fibre is trivial; hence
$\mathcal L$ is algebraically equivalent to $\mathcal O_X$.
\end{proof}

\begin{remark}\label{fga-rem-9-5-12}\dcref{9.5.12}\source{fga-explained:page-0292}\uses{}
For an arbitrary base $S$ for which $\operatorname{Pic}_{X/S}$ is locally of
finite type and represents the etale Picard functor, the relative identity
component has Lie algebra
\[
\operatorname{Lie}(\operatorname{Pic}_{X/S})
 \simeq (e^*\Omega^1_{\operatorname{Pic}_{X/S}/S})^\vee
 \simeq R^1f_*\mathcal O_X
\]
along the identity section (under the usual cohomological flatness hypotheses).
\end{remark}

\begin{corollary}\label{fga-cor-9-5-13}\dcref{9.5.13}\source{fga-explained:page-0293}\uses{fga-thm-9-4-8}
If $S=\operatorname{Spec}k$ and $\operatorname{Pic}_{X/k}$ represents the etale
Picard functor, then
\[
\dim\operatorname{Pic}^0_{X/k}\le h^1(X,\mathcal O_X).
\]
Equality holds exactly when $\operatorname{Pic}_{X/k}$ is smooth at the identity;
in that case $\operatorname{Pic}_{X/k}$ is smooth of this dimension everywhere.
\end{corollary}
\begin{proof}
The tangent space at the identity is $H^1(X,\mathcal O_X)$, so the dimension is
bounded by its dimension.  Translation identifies all components with the
identity component, and over an algebraically closed field regularity and
smoothness agree.  This proves both assertions.
\end{proof}

\begin{corollary}\label{fga-cor-9-5-14}\dcref{9.5.14}\source{fga-explained:page-0293}\uses{fga-cor-9-5-13}
If $k$ has characteristic zero, $\operatorname{Pic}_{X/k}$ is smooth of
dimension $h^1(X,\mathcal O_X)$ everywhere.
\end{corollary}
\begin{proof}
Every finite-type group scheme over a characteristic-zero field is smooth by
Cartier's theorem.  Apply Corollary~\ref{fga-cor-9-5-13}.
\end{proof}

\begin{remark}\label{fga-rem-9-5-15}\dcref{9.5.15}\source{fga-explained:page-0293}\uses{}
In positive characteristic the Picard scheme can be non-smooth even for smooth
projective surfaces.  Its tangent space is the kernel of the Bockstein operator
on $H^1(\mathcal O_X)$; the dimensions of $\operatorname{Pic}^0$, $H^1(\mathcal
O_X)$, and $H^2(\mathcal O_X)$ need not agree.  In characteristic zero, Hodge
theory and the comparison theorem impose the usual equalities of Hodge numbers.
\end{remark}

\begin{remark}[semiregularity and the characteristic system]\label{fga-rem-9-5-18}\dcref{9.5.18}\source{fga-explained:page-0294}\uses{}
Assume that $S=\operatorname{Spec}(k)$ for a field $k$, and that $X/k$ is
projective and geometrically integral.  Let $D\subset X$ be an effective
divisor with normal sheaf $\mathcal N_D$.  Write
$\delta\in\operatorname{Div}_{X/k}$ for the point representing $D$ and
$\lambda\in\operatorname{Pic}_{X/k}$ for the point representing
$\mathcal O_X(D)$.  Then
\[
T_\delta\operatorname{Div}_{X/k}=H^0(\mathcal N_D).
\]
The fundamental exact sequence
\[
0\longrightarrow\mathcal O_X\longrightarrow\mathcal O_X(D)
\longrightarrow\mathcal N_D\longrightarrow0
\]
induces
\[
\begin{aligned}
0&\longrightarrow H^0(\mathcal O_X)\longrightarrow H^0(\mathcal O_X(D))
\longrightarrow H^0(\mathcal N_D)\xrightarrow{\partial^0}H^1(\mathcal O_X)\\
&\longrightarrow H^1(\mathcal O_X(D))\xrightarrow{u}H^1(\mathcal N_D)
\xrightarrow{\partial^1}H^2(\mathcal O_X).
\end{aligned}
\]
The map $\partial^0$ is the tangent map of the Abel map
$T_\delta\operatorname{Div}_{X/k}\to T_\lambda\operatorname{Pic}_{X/k}$.
The divisor $D$ is semiregular if and only if $\partial^1$ is injective,
equivalently $u=0$.  Equivalently,
\[
\dim H^0(\mathcal N_D)=R,\qquad
R=\dim H^1(\mathcal O_X)+\dim H^0(\mathcal O_X(D))-1
 -\dim H^1(\mathcal O_X(D)).
\]
The completeness theorem for the characteristic system states that
$\operatorname{Div}_{X/k}$ is smooth of dimension $R$ at $\delta$ if and only if
$D$ is semiregular, provided $\operatorname{char}(k)=0$ or
$\operatorname{Pic}_{X/k}$ is smooth.  If $H^1(\mathcal O_X(D))=0$, the Abel map
is smooth; hence smoothness of the divisor scheme at $\delta$ is equivalent to
smoothness of the Picard scheme.  If $H^1(\mathcal N_D)=0$, the divisor scheme
is smooth at $\delta$ in every characteristic.
\end{remark}

\begin{remark}\label{fga-rem-9-5-21}\dcref{9.5.21}\source{fga-explained:page-0296}\uses{}
If $f:X\to S$ is smooth and proper in characteristic zero, Deligne's comparison
theorem makes all $R^if_*\mathcal O_X$ locally free and compatible with base
change.  Consequently the relative identity component of the Picard scheme is
smooth over each connected component of $S$; its fibre dimensions are constant
there.  If $f$ is projective Zariski-locally on $S$, the whole relative Picard
scheme is smooth, even when $S$ is nonreduced.
\end{remark}

\begin{remark}\label{fga-rem-9-5-24}\dcref{9.5.24}\source{fga-explained:page-0299}\uses{}
Let $X/S$ be an abelian scheme of relative dimension $g$.  If $X/S$ is
projective, then $\operatorname{Pic}^0_{X/S}$ is a projective abelian scheme of
dimension $g$, and the dual of $\operatorname{Pic}^0_{X/S}$ is canonically $X$.
The assertion follows fibrewise and then by proper flat descent; a normalized
Poincare sheaf realizes the duality morphism.
\end{remark}

\begin{remark}\label{fga-rem-9-5-25}\dcref{9.5.25}\source{fga-explained:page-0299}\uses{fga-prop-9-5-3,fga-thm-9-5-4}
If $X/k$ is normal, integral, and projective over an algebraically closed field,
then $\operatorname{Pic}^0_{X/k}$ is irreducible and projective.  Its reduction
is an abelian variety, and the whole identity component is an abelian variety
when the Picard scheme is reduced.
\end{remark}

\begin{remark}\label{fga-rem-9-5-26}\dcref{9.5.26}\source{fga-explained:page-0300}\uses{}
For a smooth proper family of connected curves, $J=\operatorname{Pic}^0_{X/S}$
and its dual $J^*=\operatorname{Pic}^0_{J/S}$ are projective abelian schemes.
An invertible sheaf of fibrewise degree one defines the Abel map
$X\to J$ by the graph divisor, and pullback induces the autoduality isomorphism
$J\simeq J^*$.  The construction is independent of the chosen degree-one sheaf.
\end{remark}

\begin{remark}\label{fga-rem-9-5-27}\dcref{9.5.27}\source{fga-explained:page-0300}\uses{}
For a proper flat family of curves with reducible geometric fibres, the Picard
functor is represented by an algebraic space and its separated quotient by a
separated scheme.  Raynaud's theorem assumes a normal model over a discrete
valuation ring and coprime component multiplicities; Deligne's theorem assumes
semistable (reduced nodal connected) fibres and gives a smooth quasi-projective
Picard scheme with a canonical ample line bundle.
\end{remark}

\begin{remark}\label{fga-rem-9-6-14}\dcref{9.6.14}\source{fga-explained:page-0306}\uses{fga-prop-9-6-12,fga-cor-9-4-18-3}
When $X/k$ is complete but not projective, Chow's lemma supplies a projective
surjection $X'\to X$.  The finite-type image of $\operatorname{Pic}^0_{X/k}$ in
$\operatorname{Pic}_{X'/k}$ and the numerical pullback criterion show that
$\operatorname{Pic}^0_{X/k}$ and the numerical component are still open
subschemes of finite type.  Thus the projective finiteness theorems extend to
complete schemes.
\end{remark}

\begin{remark}\label{fga-rem-9-6-19}\dcref{9.6.19}\source{fga-explained:page-0308}\uses{fga-thm-9-6-16}
If $X/S$ is proper and the Picard functor is represented, the torsion component
$\operatorname{Pic}^{\tau}_{X/S}$ is an open group subscheme of finite type.
Over a noetherian base its geometric torsion subgroup has finite bounded order;
the rank of the Neron--Severi group is finite and bounded, although it need not
be constructible as a function of the base point.
\end{remark}

\begin{remark}\label{fga-rem-9-6-22}\dcref{9.6.22}\source{fga-explained:page-0309}\uses{}
For an abelian $S$-scheme, $\operatorname{Pic}^{\tau}_{X/S}
=\operatorname{Pic}^0_{X/S}$.  The equality can be checked on geometric fibres
and then descended over $S$.
\end{remark}

\begin{remark}\label{fga-rem-9-6-24}\dcref{9.6.24}\source{fga-explained:page-0309}\uses{fga-thm-9-6-20}
For a projective flat family with integral geometric fibres of dimension $r$,
quasi-compactness of a subset of the Picard scheme can be tested by boundedness
of the Hilbert polynomial coefficients of the represented line bundles.  It is
enough to bound the leading two coefficients above and below and the next
coefficient below; equivalent criteria use the corresponding intersection
numbers with a fixed ample class.  This also implies finite type for the map
induced by a relative effective divisor.
\end{remark}

\begin{remark}\label{fga-rem-9-6-26}\dcref{9.6.26}\source{fga-explained:page-0310}\uses{}
Projectivity in the finite-type corollary cannot be weakened to properness or
local projectivity.  A reducible base curve whose components meet twice, with a
family projective near each component, gives a Picard scheme having a connected
component that is not of finite type.
\end{remark}

\begin{corollary}\label{fga-cor-9-6-27}\dcref{9.6.27}\source{fga-explained:page-0310}\uses{fga-thm-9-6-20}
If $X/S$ is projective and flat with integral geometric fibres, then for every
$n\ge0$ the power map
\[
[n]:\operatorname{Pic}_{X/S}\longrightarrow\operatorname{Pic}_{X/S},
\qquad\mathcal L\longmapsto\mathcal L^{\otimes n},
\]
is of finite type.
\end{corollary}
\begin{proof}
The identity component is of finite type.  For a connected component $U$, an
invertible sheaf representing a point of $[n]^{-1}(U)$ has Hilbert polynomial
whose relevant coefficients are bounded in terms of those of $U$ and $n$.
The numerical boundedness criterion of Theorem~9.6.20 therefore makes
$[n]^{-1}(U)$ finite type; covering the target by components proves the claim.
\end{proof}

\begin{remark}\label{fga-rem-9-6-28}\dcref{9.6.28}\source{fga-explained:page-0310}\uses{fga-cor-9-6-27}
The same finite-type assertion holds whenever $X/S$ is proper and the fppf Picard
functor is represented, even without projectivity.  It follows by the reduction
to the projective case used for the finiteness theorem and Corollary~\ref{fga-cor-9-6-27}.
\end{remark}
